% Distilled blueprint for Evans, Appendix A (Notation).

\section{Basic Notation}
\label{sec:basic-notation}

\subsection{Notation for Matrices}
\label{sec:notation-for-matrices}

\begin{definition}[Matrix Notation]
\label{def:matrix-notation-basic}
\dcref{appA:A.1-def-1}
\group{Matrix Notation}
Let $A=((a_{ij}))$ be a real $m\times n$ matrix, with $a_{ij}$ its $(i,j)$-entry; row indices may instead be written as superscripts. Then:

(i) $\operatorname{diag}(d_1,\ldots,d_n)$ denotes the diagonal matrix with entries $d_1,\ldots,d_n$.

(ii) $M^{m\times n}$ is the space of real $m\times n$ matrices, and $\mathbb{S}^{n\times n}$ is the space of real symmetric $n\times n$ matrices.

(iii) $\tr A$ is the trace of $A$.

(iv) $\det A$ is the determinant of $A$.

(v) $\cof A$ is the cofactor matrix of $A$.

(vi) $A^T$ is the transpose of $A$.

(vii) For $A=((a_{ij}))$ and $B=((b_{ij}))$ in $M^{m\times n}$,
\[
A:B=\sum_{i=1}^m\sum_{j=1}^n a_{ij}b_{ij},
\qquad
|A|=(A:A)^{1/2}
=\left(\sum_{i=1}^m\sum_{j=1}^n a_{ij}^2\right)^{1/2}.
\]
The latter is the Frobenius norm.
\end{definition}

\begin{definition}[Quadratic Forms and Related Matrix Notation]
\label{def:matrix-quadratic-form-notation}
\dcref{appA:A.1}
\group{Matrix Notation}
Let $A\in\mathbb{S}^{n\times n}$ and $x=(x_1,\ldots,x_n)\in\mathbb R^n$.

(i) The quadratic form associated with $A$ is
\[
x\cdot Ax=\sum_{i,j=1}^n a_{ij}x_ix_j.
\]

(ii) The notation $A\geq\theta I$ means
\[
x\cdot Ax\geq\theta|x|^2\qquad(x\in\mathbb R^n).
\]

(iii) For $A\in M^{m\times n}$ and $y\in\mathbb R^m$, the expression $yA$ may denote $A^Ty$.
\end{definition}

\section{Geometric Notation}
\label{sec:geometric-notation}

\begin{definition}[Geometric Notation]
\label{def:geometric-notation}
\dcref{appA:A.2}
\group{Geometric Notation}
(i) $\mathbb R^n$ is $n$-dimensional Euclidean space, with $\mathbb R=\mathbb R^1$.

(ii) $e_i=(0,\ldots,0,1,0,\ldots,0)$ is the $i$-th standard coordinate vector.

(iii) A point of $\mathbb R^n$ is written $x=(x_1,\ldots,x_n)$ and may be treated as a row or column vector according to context.

(iv)
\[
\mathbb R_+^n=\{x\in\mathbb R^n\mid x_n>0\},
\qquad
\mathbb R_+=\{x\in\mathbb R\mid x>0\}.
\]

(v) Write a point of $\mathbb R^{n+1}$ as $(x,t)=(x_1,\ldots,x_n,t)$, with $t=x_{n+1}$ interpreted as time. Also write $x=(x',x_n)$, where $x'=(x_1,\ldots,x_{n-1})\in\mathbb R^{n-1}$.

(vi) The letters $U,V,W$ denote open subsets of $\mathbb R^n$. The notation
\[
V\subset\subset U
\]
means $V\subset\overline V\subset U$ and $\overline V$ is compact; then $V$ is \textit{compactly contained} in $U$.

(vii) $\partial U$ and $\overline U=U\cup\partial U$ denote the boundary and closure of $U$.

(viii) $U_T=U\times(0,T]$.

(ix) $\Gamma_T=\overline U_T-U_T$ is the parabolic boundary of $U_T$.

(x)
\[
B^o(x,r)=\{y\in\mathbb R^n\mid |x-y|<r\}
\]
is the open ball centered at $x$ with radius $r>0$.

(xi) $B(x,r)$ is the closed ball centered at $x$ with radius $r>0$.

(xii)
\[
C(x,t,r)=\{(y,s)\in\mathbb R^n\times\mathbb R
\mid |x-y|\leq r,\ t-r^2\leq s\leq t\}
\]
is the closed cylinder with top center $(x,t)$, radius $r>0$, and height $r^2$.

(xiii)
\[
\alpha(n)=|B(0,1)|=\frac{\pi^{n/2}}{\Gamma(\frac n2+1)},
\qquad
|\partial B(0,1)|=n\alpha(n).
\]

(xiv) For $a,b\in\mathbb R^n$,
\[
a\cdot b=\sum_{i=1}^n a_ib_i,
\qquad
|a|=\left(\sum_{i=1}^n a_i^2\right)^{1/2}.
\]

(xv) $\mathbb C^n$ is $n$-dimensional complex space and $\mathbb C$ is the complex plane. For $z\in\mathbb C$, $\operatorname{Re}z$ and $\operatorname{Im}z$ are its real and imaginary parts.
\end{definition}

\section{Notation for Functions}
\label{sec:notation-for-functions}

\begin{definition}[Notation for Functions]
\label{def:function-notation-basic}
\uses{def:geometric-notation}
\dcref{appA:A.3}
\group{Function Notation}

Let $U\subset\mathbb R^n$ be open.

(i) A scalar function is written $u:U\to\mathbb R$, $u(x)=u(x_1,\ldots,x_n)$. It is \textit{smooth} when it is infinitely differentiable.

(ii) The relation $u\equiv v$ means pointwise identity, while $u:=v$ defines $u$ to equal $v$. The support of $u$ is denoted $\spt u$.

(iii)
\[
u^+=\max\{u,0\},
\qquad
u^-=-\min\{u,0\},
\qquad
u=u^+-u^-,
\qquad
|u|=u^++u^-.
\]
The sign function is
\[
\sgn(x)=
\begin{cases}
1 & x>0,\\
0 & x=0,\\
-1 & x<0.
\end{cases}
\]

(iv) A vector-valued function $\mathbf u:U\to\mathbb R^m$ is written
\[
\mathbf u(x)=(u^1(x),\ldots,u^m(x)),
\]
where $u^k$ is its $k$-th component.

(v) If $\Sigma\subset\mathbb R^n$ is a smooth $(n-1)$-surface, then $\int_\Sigma f\,dS$ denotes integration against surface measure. If $C\subset\mathbb R^n$ is a curve, then $\int_C f\,dl$ denotes integration against arclength.

(vi) Ball and sphere averages are
\[
\fint_{B(x,r)}f\,dy
=\frac{1}{\alpha(n)r^n}\int_{B(x,r)}f\,dy,
\qquad
\fint_{\partial B(x,r)}f\,dS
=\frac{1}{n\alpha(n)r^{n-1}}\int_{\partial B(x,r)}f\,dS.
\]

(vii) The indicator of $E$ is
\[
\chi_E(x)=
\begin{cases}
1 & x\in E,\\
0 & x\notin E.
\end{cases}
\]

(viii) A function $u:U\to\mathbb R$ is \textit{Lipschitz continuous} if
\[
|u(x)-u(y)|\leq C|x-y|\qquad(x,y\in U)
\]
for some $C$, and
\[
\operatorname{Lip}[u]
=\sup_{\substack{x,y\in U\\x\ne y}}
\frac{|u(x)-u(y)|}{|x-y|}.
\]

(ix) The convolution of $f$ and $g$ is denoted $f*g$.
\end{definition}

\subsection{Notation for Derivatives}
\label{sec:notation-for-derivatives}

\begin{definition}[Notation for Derivatives]
\label{def:derivative-notation-multiindex}
\dcref{appA:A.3-derivatives}
\group{Derivative Notation}
Let $u:U\to\mathbb R$ and $x\in U$.

(i) When the limit exists,
\[
\frac{\partial u}{\partial x_i}(x)
=\lim_{h\to0}\frac{u(x+he_i)-u(x)}{h}.
\]

(ii) Write $u_{x_i}=\partial u/\partial x_i$, and analogously
\[
u_{x_ix_j}=\frac{\partial^2u}{\partial x_i\partial x_j},
\qquad
u_{x_ix_jx_k}=\frac{\partial^3u}{\partial x_i\partial x_j\partial x_k}.
\]

(iii) A \textit{multiindex} is $\alpha=(\alpha_1,\ldots,\alpha_n)$ with nonnegative integer components and order
\[
|\alpha|=\alpha_1+\cdots+\alpha_n.
\]
Define
\[
D^\alpha u(x)
=\frac{\partial^{|\alpha|}u(x)}
{\partial x_1^{\alpha_1}\cdots\partial x_n^{\alpha_n}}
=\partial_1^{\alpha_1}\cdots\partial_n^{\alpha_n}u(x).
\]
For each nonnegative integer $k$,
\[
D^ku(x)=\{D^\alpha u(x)\mid |\alpha|=k\},
\qquad
|D^ku|=\left(\sum_{|\alpha|=k}|D^\alpha u|^2\right)^{1/2}.
\]
After choosing an order, $D^ku(x)$ may be regarded as a point of $\mathbb R^{n^k}$.

(iv) The first two cases are
\[
Du=(u_{x_1},\ldots,u_{x_n})
\]
for the gradient and
\[
D^2u=
\begin{pmatrix}
u_{x_1x_1}&\cdots&u_{x_1x_n}\\
\vdots&\ddots&\vdots\\
u_{x_nx_1}&\cdots&u_{x_nx_n}
\end{pmatrix}
\]
for the Hessian.

(v) The Laplacian is
\[
\Delta u=\sum_{i=1}^nu_{x_ix_i}=\tr(D^2u).
\]

(vi) A subscript on $D,D^2,\ldots$ specifies the differentiated variables. For $u=u(x,y)$ with $x\in\mathbb R^n$ and $y\in\mathbb R^m$,
\[
D_xu=(u_{x_1},\ldots,u_{x_n}),
\qquad
D_yu=(u_{y_1},\ldots,u_{y_m}).
\]
\end{definition}

\subsection{Function Spaces}
\label{sec:function-spaces-notation}

\begin{definition}[Function Spaces]
\label{def:function-space-notation}
\uses{def:derivative-notation-multiindex, def:geometric-notation}
\dcref{appA:A.3-function-spaces}
\group{Function Spaces}

Let $U\subset\mathbb R^n$ be open.

(i)
\[
\begin{aligned}
C(U)&=\{u:U\to\mathbb R\mid u\text{ is continuous}\},\\
\overline C(U)&=\{u\in C(U)\mid u\text{ is uniformly continuous}\},\\
C^k(U)&=\{u:U\to\mathbb R\mid u\text{ is }k\text{-times continuously differentiable}\},\\
\overline{C^k(U)}&=\{u\in C^k(U)\mid D^\alpha u\text{ is uniformly continuous for }|\alpha|\leq k\}.
\end{aligned}
\]
If $u\in\overline{C^k(U)}$, every $D^\alpha u$ with $|\alpha|\leq k$ extends continuously to $\overline U$.

(ii)
\[
C^\infty(U)=\bigcap_{k=0}^\infty C^k(U),
\qquad
\overline{C^\infty(U)}=\bigcap_{k=0}^\infty\overline{C^k(U)}.
\]

(iii) A subscript $c$ denotes compact support, as in $C_c(U)$ and $C_c^k(U)$.

(iv) For $1\leq p<\infty$,
\[
L^p(U)=\{u:U\to\mathbb R\mid u\text{ is measurable and }\|u\|_{L^p(U)}<\infty\},
\qquad
\|u\|_{L^p(U)}=\left(\int_U|u|^p\,dx\right)^{1/p}.
\]
Also
\[
L^\infty(U)=\{u:U\to\mathbb R\mid u\text{ is measurable and }\|u\|_{L^\infty(U)}<\infty\},
\qquad
\|u\|_{L^\infty(U)}=\operatorname*{ess\,sup}_U|u|,
\]
and
\[
L^p_{\mathrm{loc}}(U)=\{u:U\to\mathbb R\mid u\in L^p(V)\text{ for every }V\subset\subset U\}.
\]

(v)
\[
\|Du\|_{L^p(U)}=\||Du|\|_{L^p(U)},
\qquad
\|D^2u\|_{L^p(U)}=\||D^2u|\|_{L^p(U)}.
\]

(vi) $W^{k,p}(U)$ and $H^k(U)$ denote Sobolev spaces for $k=0,1,2,\ldots$ and $1\leq p\leq\infty$.

(vii) $C^{k,\beta}(U)$ and $C^{k,\beta}(\overline U)$ denote H\"older spaces for $k=0,1,2,\ldots$ and $0<\beta\leq1$.

(viii) For functions of $(x,t)\in U_T$,
\[
C_T^2(U_T)=\{u:U_T\to\mathbb R
\mid u,D_xu,D_x^2u,u_t\in C(U_T)\}.
\]
Membership entails continuity of these quantities up to the top $U\times\{t=T\}$.
\end{definition}

\section{Vector-Valued Functions}
\label{sec:vector-valued-functions}

\begin{definition}[Vector-Valued Functions]
\label{def:vector-valued-function-notation}
\uses{def:derivative-notation-multiindex, def:function-notation-basic}
\dcref{appA:A.4}
\group{Vector-Valued Functions}

Let $m>1$ and $\mathbf u=(u^1,\ldots,u^m):U\to\mathbb R^m$.

(i) For every multiindex $\alpha$,
\[
D^\alpha\mathbf u=(D^\alpha u^1,\ldots,D^\alpha u^m),
\qquad
D^k\mathbf u=\{D^\alpha\mathbf u\mid |\alpha|=k\},
\]
and
\[
|D^k\mathbf u|
=\left(\sum_{|\alpha|=k}|D^\alpha\mathbf u|^2\right)^{1/2}.
\]

(ii) The gradient matrix is
\[
D\mathbf u=
\begin{pmatrix}
\partial u^1/\partial x_1&\cdots&\partial u^1/\partial x_n\\
\vdots&\ddots&\vdots\\
\partial u^m/\partial x_1&\cdots&\partial u^m/\partial x_n
\end{pmatrix}_{m\times n}.
\]

(iii) If $m=n$, then
\[
\Div\mathbf u
=\tr(D\mathbf u)
=\sum_{i=1}^n u^i_{x_i}.
\]

(iv) The spaces $C(U;\mathbb R^m)$, $L^p(U;\mathbb R^m)$, and analogous spaces consist of maps whose components belong to the corresponding scalar space.
\end{definition}

\begin{remark}[Sub- and Superscripts for Boldface Mappings]
\label{rem:boldface-convention-sub-superscripts}
\uses{def:vector-valued-function-notation}
\dcref{appA:A.4-remark}
\group{Notation Conventions}

Mappings into $\mathbb R^m$ for $m>1$, or into Banach or Hilbert spaces, are boldface and have superscripted components. Points $x=(x_1,\ldots,x_n)\in\mathbb R^n$ are not boldface and have subscripted components. Matrix-valued mappings are boldface; their component indices may be superscripts or mixed sub- and superscripts.
\end{remark}

\section{Notation for Estimates}
\label{sec:notation-for-estimates}

\begin{remark}[Constants]
\label{rem:constant-notation-convention}
\dcref{appA:A.5-constants}
\group{Asymptotic Notation}
The symbol $C$ denotes a constant computable from the stated data. Its value may change between lines as fixed numerical factors are absorbed into it.
\end{remark}

\begin{definition}[Big-Oh and Little-Oh Notation]
\label{def:big-oh-little-oh-notation}
\dcref{appA:A.5-definitions}
\group{Asymptotic Notation}
(i) The relation
\[
f=O(g)\quad\text{as }x\to x_0
\]
means that some $C$ satisfies
\[
|f(x)|\leq C|g(x)|
\]
for all $x$ sufficiently close to $x_0$.

(ii) The relation
\[
f=o(g)\quad\text{as }x\to x_0
\]
means
\[
\lim_{x\to x_0}\frac{|f(x)|}{|g(x)|}=0.
\]
\end{definition}

\begin{remark}[Implicit Limits in $O$/$o$ Notation]
\label{rem:big-oh-little-oh-implicit-limit}
\uses{def:big-oh-little-oh-notation}
\dcref{appA:A.5-remark}
\group{Asymptotic Notation}

The expressions $O(g)$ and $o(g)$ require a limiting regime, such as $x\to x_0$, even when it is left implicit by context.
\end{remark}

\section{Some Comments about Notation}
\label{sec:comments-about-notation}

\begin{remark}[Why ``$Du$'' Instead of ``$\nabla u$'']
\label{rem:gradient-notation-choice-du}
\uses{def:derivative-notation-multiindex}
\dcref{appA:A.6-i}
\group{Notation Conventions}

The gradient and Hessian are denoted $Du$ and $D^2u$. This keeps $D^2u$ distinct from the Laplacian and makes the multiindex notation $D^\alpha u$ consistent.
\end{remark}

\begin{remark}[Why ``$U$'' Instead of ``$\Omega$'']
\label{rem:domain-notation-choice-u}
\uses{def:geometric-notation}
\dcref{appA:A.6-ii}
\group{Notation Conventions}

Open PDE domains are normally denoted $U$, leaving $V$ and $W$ for subdomains. The symbol $\Omega$ is reserved for probability spaces, including probabilistic representations of PDE \cite{Freidlin1985}.
\end{remark}
